% Nová příprava: 2026-10-28; 1. ročník SPŠ, etapa 13–24.
% Rozpis: 1_rocnik/tematicke_plany/Rozpis_vyuky_MAT_1_rocnik.tex, hodina 23.
% Původ: vlastní formulace výkladu a vlastní úlohy; metodická inspirace uvedenými zdroji.
% Kontrola, obsahové porovnání zdrojů a časy: planovani/DOKONCENI_PRIPRAV_13_24.md.
% Učebnice: ucebnice/03_Planimetrie/Planimetrie.tex, Trojúhelník a značení;
% ucebnice/04_Goniometrie_a_trigonometrie/Goniometrické funkce.tex,
% Tabulkové hodnoty (pouze ostré úhly) a praktické použití inverzních funkcí.
% Realisticky.cz: externi_zdroje/realisticky/01_Zakladni_poznatky/01_Opakovani/010119_Podobnosti_trojuhelniku_goniometricke_funkce/010119_Lekce.pdf
% Realisticky.cz: externi_zdroje/realisticky/01_Zakladni_poznatky/01_Opakovani/010120_Pouziti_goniometrickych_funkci_ve_fyzice/010120_Lekce.pdf
\documentclass[11pt,a4paper]{article}
% Spolecny styl pro ucitelske pripravy MAT (Work/MAT)
% Verze 1.1, 2026-10-23
\NeedsTeXFormat{LaTeX2e}
\ProvidesPackage{priprava}[2026/10/23 v1.1 Sablona ucitelskych priprav MAT]

% Zakladni typografie a cestina
\RequirePackage[T1]{fontenc}
\RequirePackage[utf8]{inputenc}
\RequirePackage[czech]{babel}
\RequirePackage{lmodern}
\RequirePackage{microtype}
\RequirePackage{amsmath,amssymb}
\RequirePackage{xcolor}
\RequirePackage{enumitem}
\RequirePackage[most]{tcolorbox}
\RequirePackage{tabularx,booktabs,array}
\RequirePackage{fancyhdr}
\RequirePackage{needspace}
\RequirePackage{geometry}
\RequirePackage[hidelinks]{hyperref}

\geometry{a4paper,top=17mm,bottom=18mm,left=23mm,right=23mm}
\definecolor{notegray}{gray}{0.39}
\setlength{\parindent}{0pt}
\setlength{\parskip}{2.5pt plus 1pt minus 0.5pt}
\setlist[itemize]{topsep=2pt,itemsep=1pt,parsep=0pt,leftmargin=1.45em}
\setlist[enumerate]{topsep=2pt,itemsep=2pt,parsep=0pt,leftmargin=1.65em}
\setcounter{secnumdepth}{0}
\renewcommand{\arraystretch}{1.12}

% Zahlavi: vlevo rocnik | globalni tema, vpravo cislo: nazev hodiny
% Zapati zustava zcela prazdne.
\pagestyle{fancy}
\fancyhf{}
\renewcommand{\headrulewidth}{0.35pt}
\renewcommand{\footrulewidth}{0pt}
\setlength{\headheight}{14pt}
\newcommand{\nastavPripravu}[4]{%
  \fancyhead[L]{\small #1\enspace\textbar\enspace #2}%
  \fancyhead[R]{\small #3: #4}%
}

% Minimalisticke cernobile boxy. Kazda definice zvlast.
\tcbset{commonlesson/.style={enhanced,breakable=false,colback=white,colframe=black,
  colbacktitle=white,coltitle=black,boxrule=0.45pt,
  arc=0.8mm,outer arc=0.8mm,left=3mm,right=3mm,top=1.6mm,bottom=1.6mm,
  toptitle=1mm,bottomtitle=0.7mm,fonttitle=\bfseries,
  attach boxed title to top left={xshift=2.5mm,yshift=-1.15mm},
  boxed title style={colback=white,colframe=white,boxrule=0pt,left=1mm,right=1mm}}}
\newtcolorbox{zadaniBox}[1][]{commonlesson,title={Zadání},#1}
\newtcolorbox{prikladBox}[1][]{commonlesson,title={Příklad},#1}
\newtcolorbox{definiceBox}[1][]{commonlesson,title={Definice},#1}

% Spolecny box pro zadani a jeho reseni, bez mezery mezi castmi.
% Pouziti:
%   \begin{prikladReseniBox}
%     Zadani ulohy.
%     \tcblower
%     \begin{align*} ... \end{align*}
%   \end{prikladReseniBox}
% Spodni cast ma jemny sedy podklad a delici linku; lze v ni sazet
% vice radku matematiky. Obe casti maji jen jeden spolecny ramecek.
\tcbset{solutionlesson/.style={
  commonlesson,
  enhanced jigsaw,
  breakable,
  colbacklower=black!3,
  segmentation style={solid,draw=black!25,line width=0.35pt},
  fontlower=\small\color{notegray},
  before lower={\noindent\textbf{Řešení:}\enspace},
  leftlower=3mm,rightlower=3mm
}}
\newtcolorbox{prikladReseniBox}[1][]{solutionlesson,title={Příklad},#1}
\newtcolorbox{zadaniReseniBox}[1][]{solutionlesson,title={Zadání},#1}

% Pomucky pro autora pripravy: sedive poznamky; \reseni pro starsi pripravy
\newcommand{\poznamka}[1]{\par\begingroup\small\color{notegray}\itshape
  \textnormal{\textbf{Poznámka pro učitele.}}\enspace #1\par\endgroup}
\newcommand{\reseni}[1]{\par\begingroup\small\color{notegray}
  \textbf{Očekávané řešení:}\enspace #1\par\endgroup}
\newcommand{\blok}[1]{\needspace{5\baselineskip}\subsection*{#1}}
\newcommand{\vyrok}[1]{\uv{#1}}

\usepackage{tikz}
\usetikzlibrary{patterns,calc,angles,quotes,babel}
\nastavPripravu{1. ročník}{Goniometrie}{23}{Sinus, kosinus a tangens}
\begin{document}
\blok{Stejný úhel, stejný poměr}
Zvětšíme-li pravoúhlý trojúhelník dvakrát, všechny jeho strany se zdvojnásobí,
ale poměr dvou stran zůstane stejný. Dva pravoúhlé trojúhelníky se stejným
ostrým úhlem jsou podobné podle uu: shodují se v tomto úhlu i v pravém úhlu.
Poměry tedy závisejí na ostrém úhlu, ne na velikosti náčrtu.

\begin{definiceBox}
\textbf{Podobné trojúhelníky} mají stejné odpovídající úhly a poměry
odpovídajících délek stran jsou shodné.
\end{definiceBox}
\begin{definiceBox}
\textbf{Protilehlá odvěsna k ostrému úhlu} leží proti jeho vrcholu.
\end{definiceBox}
\begin{definiceBox}
\textbf{Přilehlá odvěsna k ostrému úhlu} je odvěsna, která tvoří jedno jeho rameno.
Přepona také tvoří rameno úhlu, ale přilehlou odvěsnou není.
\end{definiceBox}
\begin{center}
\begin{tikzpicture}[scale=.7]
\coordinate (A) at (0,0); \coordinate (B) at (4,3); \coordinate (C) at (4,0);
\draw[thick] (A)--node[below]{$b$}(C)--node[right]{$a$}(B)--node[above,sloped]{$c$}(A);
\node[below left] at (A) {$A$}; \node[above] at (B) {$B$}; \node[below right] at (C) {$C$};
\pic[draw,angle radius=3mm]{right angle=B--C--A};
\pic[draw,angle radius=9mm,"$\alpha$",angle eccentricity=1.5]{angle=C--A--B};
\end{tikzpicture}
\end{center}
Zde $\gamma=90^\circ$: vzhledem k $\alpha$ je $a$ protilehlá odvěsna,
$b$ přilehlá a $c$ přepona. Vzhledem k druhému ostrému úhlu se role odvěsen vymění.
Otočení obrázku toto pravidlo nemění.

\begin{definiceBox}
Pro $0^\circ<\alpha<90^\circ$ je \textbf{sinus}
$\displaystyle\sin\alpha=\frac{\text{protilehlá odvěsna}}{\text{přepona}}$.
\end{definiceBox}
\begin{definiceBox}
Pro $0^\circ<\alpha<90^\circ$ je \textbf{kosinus}
$\displaystyle\cos\alpha=\frac{\text{přilehlá odvěsna}}{\text{přepona}}$.
\end{definiceBox}
\begin{definiceBox}
Pro $0^\circ<\alpha<90^\circ$ je \textbf{tangens}
$\displaystyle\tan\alpha=\frac{\text{protilehlá odvěsna}}{\text{přilehlá odvěsna}}$.
Setkáme se také se značkou $\operatorname{tg}$.
\end{definiceBox}
Všechny délky jsou kladné a ve stejných jednotkách, takže poměry nemají jednotku.
Sinus i kosinus ostrého úhlu jsou mezi $0$ a $1$, tangens je kladný
a může být větší než $1$. Podobnost z úvodu zaručuje, že definice nezávisí
na zvoleném trojúhelníku.

\begin{prikladReseniBox}[title={P1: přesné poměry},breakable=false]
Pro $\gamma=90^\circ$, $a=5\,\mathrm{cm}$, $b=12\,\mathrm{cm}$,
$c=13\,\mathrm{cm}$ urči sinus, kosinus a tangens obou ostrých úhlů.
\tcblower
$\sin\alpha=\frac5{13}$, $\cos\alpha=\frac{12}{13}$,
$\tan\alpha=\frac5{12}$.
Pro $\beta$ se role odvěsen vymění:
$\sin\beta=\frac{12}{13}$, $\cos\beta=\frac5{13}$,
$\tan\beta=\frac{12}5$. Přepona zůstává stejná.
\end{prikladReseniBox}

\blok{Z poměru ke konkrétní délce a úhlu}
Na kalkulačce nastav stupně, obvykle \texttt{DEG}.
Kontrolní hodnota je $\sin30^\circ=0{,}5$.
Při hledání délky nejprve napiš poměr a vyjádři neznámou, až pak dosaď.
Zaokrouhluj teprve konečné délky na setiny a úhly na desetiny stupně.

\begin{prikladReseniBox}[title={P2: přepona a ostrý úhel},breakable=false]
Pro $\gamma=90^\circ$, $c=18\,\mathrm{cm}$ a $\alpha=35^\circ$ urči $a,b,\beta$.
\tcblower
\[
\frac ac=\sin\alpha\Rightarrow a=c\sin\alpha
=18\sin35^\circ\approx10{,}32\,\mathrm{cm},
\]
\[
\frac bc=\cos\alpha\Rightarrow b=c\cos\alpha
=18\cos35^\circ\approx14{,}74\,\mathrm{cm}.
\]
$\beta=55^\circ$. Obě odvěsny jsou kratší než $18\,\mathrm{cm}$.
Z nezaokrouhlených hodnot také $a^2+b^2=c^2$.
\end{prikladReseniBox}

\begin{prikladReseniBox}[title={P3: hledáme úhel z poměru},breakable=false]
Pravoúhlý trojúhelník má vzhledem k $\alpha$ protilehlou odvěsnu
$7\,\mathrm{cm}$ a přilehlou $11\,\mathrm{cm}$. Urči oba ostré úhly.
\tcblower
$\tan\alpha=\frac7{11}$. Pomocí inverzní funkce dostaneme
$\alpha=\arctan\frac7{11}\approx32{,}5^\circ$ a druhý úhel přibližně $57{,}5^\circ$.
Kontrola: menší odvěsna leží proti menšímu ostrému úhlu a součet je $90^\circ$.
Tlačítka $\tan^{-1}$ či $\sin^{-1}$ označují hledání úhlu, nikoli převrácení
hodnoty tangens nebo sinus. V tomto kroku hledáme pouze ostrý úhel.
\end{prikladReseniBox}
\poznamka{Jádro: názorná podobnost, definice, P1–P3 a S1 a), b), d), e),
asi 38 min práce, organizace a závěr 7 min. S1 c) je další nabídka.
Definice neopisuj třikrát z tabule celé: nech žáky ukázat příslušné strany
a zapsat poměr. Jednotkovou kružnici, radiány ani obecné goniometrické rovnice nezavádíme.}

\blok{Samostatné procvičení}
\begin{zadaniReseniBox}[title={S1: pojmenuj strany před dosazením},breakable=false]
V a)–c) je $\gamma=90^\circ$.
\begin{enumerate}[label=\alph*)]
\item Pro $a=8$, $b=6$, $c=10\,\mathrm{cm}$ urči tři poměry pro $\alpha$ a $\beta$.
\item Pro $c=9\,\mathrm{cm}$ a $\alpha=42^\circ$ urči obě odvěsny.
\item Pro $b=12\,\mathrm{cm}$ a $\alpha=28^\circ$ urči $a$ a $c$.
\item Urči ostrý úhel, je-li jeho sinus $0{,}4$. Změní se výsledek, je-li $0{,}4$ jeho kosinus?
\item Může být sinus ostrého úhlu roven $1{,}2$? Vysvětli pomocí stran.
\end{enumerate}
\tcblower
a) Pro $\alpha$: $\frac45,\frac35,\frac43$; pro $\beta$:
$\frac35,\frac45,\frac34$, vždy v pořadí sinus, kosinus, tangens.
b) $a=9\sin42^\circ\approx6{,}02\,\mathrm{cm}$,
$b=9\cos42^\circ\approx6{,}69\,\mathrm{cm}$.
c) $a=12\tan28^\circ\approx6{,}38\,\mathrm{cm}$,
$c=\frac{12}{\cos28^\circ}\approx13{,}59\,\mathrm{cm}$.
d) $\arcsin0{,}4\approx23{,}6^\circ$, kdežto $\arccos0{,}4\approx66{,}4^\circ$.
Název funkce je součást údaje.
e) Ne, protilehlá odvěsna je kratší než přepona, proto jejich poměr je menší než $1$.
\end{zadaniReseniBox}

\newpage
\blok{Volitelná rezerva}
\begin{zadaniReseniBox}[title={R1: odvození zvláštních hodnot, 12–18 minut},breakable=false]
Odvoď hodnoty pro $45^\circ$ z rovnoramenného pravoúhlého trojúhelníku
s odvěsnami $1$ a pro $30^\circ,60^\circ$ z poloviny rovnostranného trojúhelníku o straně $2$.
\tcblower
První přepona je $\sqrt2$. Ve druhém případě má polovina základny délku $1$,
výška je $\sqrt{2^2-1^2}=\sqrt3$; proti úhlu $30^\circ$ leží kratší odvěsna.
\[
\begin{array}{c|ccc}
\alpha&30^\circ&45^\circ&60^\circ\\\hline
\sin\alpha&\frac12&\frac1{\sqrt2}&\frac{\sqrt3}{2}\\[1mm]
\cos\alpha&\frac{\sqrt3}{2}&\frac1{\sqrt2}&\frac12\\[1mm]
\tan\alpha&\frac1{\sqrt3}&1&\sqrt3
\end{array}
\]
Podoby s odmocninou ve jmenovateli zde ponecháme; usměrňování budeme učit později.
\end{zadaniReseniBox}
\begin{zadaniReseniBox}[title={R2: pravý úhel u jiného vrcholu},breakable=false]
Je-li $\alpha=90^\circ$, $\beta=40^\circ$ a $a=20\,\mathrm{cm}$, urči $b,c,\gamma$.
\tcblower
Přepona je $a$. Pro úhel $\beta$ je $b$ protilehlá a $c$ přilehlá odvěsna.
$b=20\sin40^\circ\approx12{,}86\,\mathrm{cm}$,
$c=20\cos40^\circ\approx15{,}32\,\mathrm{cm}$, $\gamma=50^\circ$.
Nelze mechanicky použít $\sin\beta=b/c$, protože $c$ není přepona.
\end{zadaniReseniBox}
\poznamka{R2 je další rezerva na 5–7 min, nikoli povinné pokračování R1.}

\blok{Domácí procvičení}
\begin{zadaniReseniBox}[title={D1: 15–20 minut; kalkulačka ve stupních}]
V a)–d) je $\gamma=90^\circ$.
\begin{enumerate}[label=\alph*)]
\item Pro $a=12$, $b=5$, $c=13\,\mathrm{cm}$ zapiš tři poměry pro $\alpha$.
\item Pro $c=16\,\mathrm{cm}$ a $\alpha=25^\circ$ urči $a,b,\beta$.
\item Pro $b=9\,\mathrm{cm}$ a $\alpha=38^\circ$ urči $a,c$.
\item Pro $c=15\,\mathrm{cm}$ a $a=9\,\mathrm{cm}$ urči oba ostré úhly.
\item Jak ověříš režim kalkulačky před výpočtem?
\end{enumerate}
\tcblower
a) $\sin\alpha=\frac{12}{13}$, $\cos\alpha=\frac5{13}$, $\tan\alpha=\frac{12}5$.
b) $a=16\sin25^\circ\approx6{,}76\,\mathrm{cm}$,
$b=16\cos25^\circ\approx14{,}50\,\mathrm{cm}$, $\beta=65^\circ$.
c) $a=9\tan38^\circ\approx7{,}03\,\mathrm{cm}$,
$c=\frac9{\cos38^\circ}\approx11{,}42\,\mathrm{cm}$.
d) $\alpha=\arcsin\frac9{15}\approx36{,}9^\circ$, $\beta\approx53{,}1^\circ$.
e) Zkontroluji \texttt{DEG} a zkusím $\sin30^\circ=0{,}5$.
\end{zadaniReseniBox}
\end{document}
